% =====================================================================
\chapter{Problem Identification and Innovation}
\label{sec:problem}
% =====================================================================

\section{The problem: nobody checks the checker}

When my FSc Part-I results came out, my Physics marks were much lower than I expected. I knew something was wrong.
I had not just answered the paper; I had attempted three extra questions, and I had written every answer neatly so the
examiner could follow each step. So I applied for rechecking.

When I finally saw my paper, I was shocked. Almost every answer had a tick on it. Nothing was crossed out in red. And
still, I had been given about half the marks. I pointed this out to the official at the board office. He said,
``This is recounting, not rechecking.'' He meant that they only add up the marks again. Nobody looks at whether the
marking itself was right. My marks stayed the same.

That was the first real shock of my education. It also made me notice something I had heard about for years: in
government schools, teachers who mark board papers are paid for each paper they check, so checking more papers
means earning more. Fast marking is rewarded. Careful marking is not.

And this is not a small mistake. In our system, a few marks decide almost everything: which college a student gets
into, whether they make the merit list for medicine or engineering, which path their life takes. One careless checking,
done in a hurry by someone trying to finish more papers, can push a student away from the direction they wanted to
follow. Often they never even find out why. That is not a small damage. It can change the course of a whole life.

Since then, one question has stayed with me. If an examiner makes a mistake, who notices? In most cases, nobody does.
In our board exams, nobody checks the checker.

\begin{keybox}[The problem in one line]
In Pakistan's board exams, millions of handwritten scripts are marked by examiners who are paid per script, and most
scripts are never seen by a second person. When a mark is wrong, the student usually has no real way to get it corrected.
\end{keybox}

\section{Why it is worth solving}
\begin{itemize}
  \item \textbf{Scale.} In 2026, 1,348,352 candidates sat Punjab's Class 10 exam and 1,395,575 sat Class 9 across the nine
        Punjab boards \citep{propk2026matric,propk2026class9}. At roughly 8 theory papers each, that is about
        \textbf{20 million answer scripts a year} for matric in Punjab alone (my estimate).
  \item \textbf{Stakes.} These marks decide merit lists, college admissions and MDCAT/ECAT aggregates. In one reported
        case, a candidate's mark went from 2/75 to 61/75 after rechecking \citep{tribune2014fiasco} (a 2014 report).
  \item \textbf{Incentives reward speed.} Under Punjab's 2023 pay revision, matric examiners receive \rs45 per paper if they
        check more than 50 papers and \rs35 if they check fewer \citep{propk2023fee}. The rate goes up for marking more.
  \item \textbf{Errors are documented, not hypothetical.} After the SSC 2024 exams, BISE Rawalpindi penalised 435 marking
        staff for wrong marking, including 66 head examiners and 210 sub-examiners \citep{dawn2024examiners}. In December
        2025, a Class 9 student found an unmarked question and answers ``ticked and under-credited without justification''
        \citep{tribune2025receiving}. That is the same pattern as my own paper.
  \item \textbf{Results are falling.} Rawalpindi board had the lowest pass rate in Punjab in 2026: 37.56\% in Class 9
        \citep{propk2026class9}. Marking is not the only reason, but when results fall, students deserve to know their
        marks are at least correct.
  \item \textbf{The system cannot measure its own fairness.} I could not find any board publishing how many rechecking
        requests it receives or how many marks change after review \citep{research01}.
\end{itemize}

\section{How the problem is addressed today}

\begin{table}[H]
\centering
\caption{Current ways of checking marking quality, in Pakistan and abroad.}
\label{tab:current}
\begin{tabularx}{\linewidth}{@{}P{4.6cm}L@{}}
\toprule
\textbf{Mechanism} & \textbf{What it does} \\
\midrule
Manual marking centres (Punjab) & Head examiner and sub-examiners; the head examiner rechecks a sample of about 17--20\% of
  scripts \citep{bisemultan,ilmkidunya2026} \\
Rechecking (Punjab, until 2025) & Only verified totals and unmarked questions, with no review of marking accuracy
  \citep{techjuice2025remark} \\
Re-marking (BISE Lahore, announced Oct 2025) & Real re-evaluation, pending cabinet approval \citep{propk2025remark} \\
E-marking (Punjab, from 2026, pilots first) & Scanned scripts, secure examiner IDs, face verification, time limits; the
  system adds up totals \citep{dawn2025digitisation,tribune2026online} \\
E-marking (BISE Peshawar, 2023) & First examiner marks; a second randomly reviews 20\% of marked papers; the head
  examiner intervenes \citep{dawn2023peshawar} \\
E-marking (FBISE) & About 10\% of papers go to the head examiner; a head examiner said their checking role became more limited
  under e-marking \citep{dawn2023fbise} \\
International (Cambridge, AQA) & ``Seeding'': pre-marked scripts hidden among examiners' work to test them
  \citep{cambridgeqa,ofqual2018} \\
International research (AQA, Ofqual) & Trials of AI as a \emph{quality check} of human marking, not as a marker
  \citep{tes2024aqa,kcl_aqa,ofqual2026} \\
\bottomrule
\end{tabularx}
\end{table}

\section{Limitations of existing solutions}
\begin{enumerate}
  \item \textbf{Checks are random samples.} 80--90\% of scripts are never looked at by a second person. AQA's own research
        found that sampling-based quality assurance has little potential to improve marking reliability directly: it
        catches bad examiners for retraining, but does not fix the wrong mark on a particular student's paper
        \citep{pinotdemoira}.
  \item \textbf{E-marking fixes arithmetic, not judgement.} The system adds up totals itself \citep{ofqual2018}, but nothing
        public says how Punjab's version checks whether a correct answer was under-credited.
  \item \textbf{The transition is long.} Punjab expects two to three years before all exams move to e-marking
        \citep{pakobserver2026}. Until then, most subjects are marked by hand on paper.
  \item \textbf{Rechecking is reactive and costly for families.} The student has to suspect an error, pay, travel and argue.
        A wrong mark on a paper nobody suspects is never found.
  \item \textbf{UK AI work does not transfer directly.} The idea of AI checking human markers is not new: AQA has trialled
        it, and its chief executive said AI could ``mark markers'' very soon \citep{tes2024aqa,schoolsweek2026}. But that
        work assumes mature on-screen systems and clean data. Pakistani scripts are handwritten, bilingual (Urdu and
        English), and often still on paper \dec{D006}.
\end{enumerate}

\section{My approach}
I do not replace examiners, and I do not let AI give marks. Regulators currently describe AI as promising for quality
assurance but not ready to mark high-stakes exams \citep{ofqual2026}. So I propose Parakh, an \textbf{audit layer} that tells the
existing reviewers \emph{where to look} \dec{D004,D006}:
\begin{enumerate}
  \item \textbf{Consistency rules (no AI judgement).} Check the examiner's own marking on the script: answers ticked but
        marks cut, unmarked questions, totals that do not add up, marks above the maximum.
  \item \textbf{AI second reader.} A vision-language model reads each answer and estimates a mark from the answer key.
        A big gap between its estimate and the examiner's mark raises the script's risk score. It never sets the mark.
  \item \textbf{Risk-ranked review instead of random review.} Boards already pay for re-checking a share of scripts. The
        same effort goes to the highest-risk scripts first.
  \item \textbf{On-premise.} Scripts are secret and there are millions of them, so models run on the board's own machines.
\end{enumerate}

\begin{keybox}[Hypotheses tested in the prototype]
\textbf{H1:} With a fixed review budget (20\% of scripts re-checked by a human), risk-ranked review catches more marking
errors than random review. \\
\textbf{H2:} Auditing a script by machine costs much less than one human re-check. \dec{D007}
\end{keybox}
